In this report, we presented two structured products on \gls{jpm}, a \glsreset{pppn}\gls{pppn} and an Airbag Note, both built from a zero-coupon bond and vanilla European options.
Table~\ref{tab:structured-products-comparison} compares them.

\begin{table}[H]
    \centering
    \caption{Comparison of the two structured products.}
    \label{tab:structured-products-comparison}
    \begin{tabular}{@{}>{\raggedright\arraybackslash}p{3.5cm}>{\raggedright\arraybackslash}p{6cm}>{\raggedright\arraybackslash}p{6cm}@{}}
    \toprule
    \textbf{Aspect} & \textbf{PPPN} & \textbf{Airbag Note} \\
    \midrule
    Structure &
    90\% capital protection. Pays a premium if the stock moves out of the corridor $[270, 350]$ in either direction. &
    Full capital protection down to the airbag level $A = 256.33$. Participates partially ($p = 0.690$) in gains. \\[2mm]

    Replication &
    Zero-coupon bond worth $0.9\,N e^{-rT}$, long $EP(270,T)$ and long $EC(350,T)$. &
    Zero-coupon bond worth $N e^{-rT}$, long $EC(S_0,T)$ and short $EP(A,T)$. \\[2mm]

    Issuer margin & $0.5\%$ & $0.5\%$ \\[2mm]

    Expected growth (Monte Carlo) &
    $\hat{\mathbb{E}}[\text{Payoff}(S_T)]/N = 1.198$ &
    $\hat{\mathbb{E}}[\varphi_{p,A}(S_T)] = 1.234$ \\
    \bottomrule
    \end{tabular}
\end{table}

The \gls{pppn} combines partial capital protection with exposure to large moves in either direction.
The Airbag Note instead focuses on downside protection: the investor keeps the full capital as long as \gls{jpm} stays above the airbag level and participates in its gains, while moderate losses are absorbed.
Both products are therefore aimed at risk-conscious rather than purely risk-seeking investors, that is, at investors who want their capital protected to some degree but still want to benefit from rising prices.

Monte Carlo simulations based on the historical drift of \gls{jpm} point to favourable expected performance for both products.
For the \gls{pppn}, the present value of the expected premium is more than twice the cost of the two options.
For the Airbag Note, the expected payoff factor is above one, and a loss occurs in fewer than one in ten simulated scenarios.
These estimates rely on \gls{jpm}'s strong performance over the past five years continuing, so they should be read as indicative rather than as a forecast.
